%! TEX program = pdflatex
% =====================================================================
% Internal document: full statements and detailed proofs of every
% mathematical claim in the manuscript, with verification notes.
% Machine companion: analysis/verify_proofs.py (sympy + Monte Carlo).
% Not part of any submission.
% =====================================================================
\documentclass[11pt,a4paper]{article}
\usepackage[utf8]{inputenc}
\usepackage[T1]{fontenc}
\usepackage{lmodern}
\usepackage[margin=1.05in]{geometry}
\usepackage{amsmath,amssymb,amsthm,mathtools}
\usepackage{xcolor}
\usepackage[hidelinks]{hyperref}
\usepackage{enumitem}
\usepackage{microtype}
\usepackage{booktabs}

\theoremstyle{plain}
\newtheorem{claim}{Claim}
\theoremstyle{remark}
\newtheorem*{verification}{Verification status}

\newcommand{\E}{\mathbb{E}}
\newcommand{\PASS}{\textcolor{green!50!black}{\textbf{VERIFIED}}}
\newcommand{\FIXED}{\textcolor{orange!80!black}{\textbf{ERROR FOUND \& FIXED}}}

\title{\textbf{Proof Verification Dossier}\\[4pt]
\large Bounds and Estimators for the Civilian-to-Combatant Death Ratio\\[2pt]
\normalsize Internal document for manual checking --- not for submission.\\
Machine companions: \texttt{analysis/verify\_proofs.py} (sympy $+$ Monte Carlo, all checks pass)\\
and \texttt{formal/casualty\_proofs} (Lean~4 $+$ mathlib, kernel-checked, no \texttt{sorry}s).}
\author{}
\date{\today}

\begin{document}
\maketitle

\noindent\textbf{How to read this document.} Each claim is stated exactly as in the manuscript, followed by a proof written in more detail than the paper's appendix, an explicit list of the assumptions used at each step, and a verification note describing the symbolic (sympy) or Monte-Carlo check performed. The algebraic claims (1--4, the $M$-monotonicity step of Claim~7, and the Section~6 arithmetic over exact rationals) are additionally machine-checked in Lean~4 with mathlib (\texttt{formal/casualty\_proofs}); the Lean development compiles with no unproven statements. Three errors were found in the first verification round and corrected in the manuscript; a second round of independent external review of this dossier surfaced further defects (most materially in the former Proposition~5.1), also now corrected in both dossier and manuscript. All are flagged inline and catalogued in the final two sections.

\tableofcontents

%======================================================================
\section{Setup and standing conventions}

Population of size $N$; deaths $D$ in the window $[t_0,t_1]$. Classes: $W$ (women of all ages $+$ males under 18) with population share $w$, and $\mathrm{AM}$ (males 18$+$) with population share $a$. \emph{Convention (revised in this version):} $\mathrm{AM}$ is \emph{all} males 18$+$, so
\[
w + a = 1 \quad\text{exactly.}
\]
(The prior draft used $a=$ males 18--60 $=0.255$, leaving men over 60 unassigned; this left $w+a=0.988$ and made $a=1-w$ substitutions silently wrong. Now $a=0.267$.)

Combatant stock $M$, $f = M/N$, all combatants in $\mathrm{AM}$ (slack for exceptions handled in Claim~\ref{cl:remark}). Estimand $q = D_M/D$, the combatant share of the dead; $\omega$ = share of the \emph{dead} in class $W$.

\paragraph{The hazard model.} All identification results use one structural model: conditional on being a civilian, a member of class $W$ dies with per-capita hazard $\lambda$ and a member of civilian-$\mathrm{AM}$ with hazard $\mu\lambda$, $\mu\ge 1$. Combatant deaths are governed separately (their total is $qD$ by definition). No assumption is made about $\lambda$ itself --- it cancels.

%======================================================================
\section{Claim 1: benchmark identification (paper eq.~2, Assumption 3.1)}

\begin{claim}[Point identification under exchangeable collateral]\label{cl:ident}
If all civilians face equal per-capita hazard ($\mu=1$), then
\[
q = 1-\frac{\omega(1-f)}{w}.
\]
\end{claim}

\begin{proof}
Civilian population masses: class $W$ has $wN$ civilians (no combatants in $W$); civilian $\mathrm{AM}$ has $(a-f)N$. With equal hazard $\lambda$, expected civilian deaths in $W$ and civilian-$\mathrm{AM}$ are $\lambda wN$ and $\lambda(a-f)N$, so the fraction of \emph{civilian} deaths landing in $W$ is
\[
\frac{wN}{wN+(a-f)N} = \frac{w}{w+a-f} = \frac{w}{1-f},
\]
where the last step uses $w+a=1$. Combatant deaths (share $q$ of all deaths) land entirely in $\mathrm{AM}$. Hence the share of \emph{all} deaths in $W$ is
\[
\omega = (1-q)\cdot\frac{w}{1-f},
\]
and solving for $q$ gives the display. Note this is exact in expectation (a first-moment identity), not an approximation; finite-sample noise is handled by the delta method (Claim~\ref{cl:delta}).
\end{proof}

\begin{verification}
\PASS. Sympy: solving $\omega=(1-q)w\lambda/(w\lambda+(a-f)\mu\lambda)$ for $q$ and substituting $\mu=1$, $a=1-w$ reproduces the display exactly ($\lambda$ cancels symbolically). Degenerate cases: $\omega=w$, $f=0$ $\Rightarrow$ $q=0$; $\omega=0$ $\Rightarrow$ $q=1$. \emph{Note the subtlety}: the formula requires $a=1-w$; with the old convention $a=0.255\neq 1-w$ the two displayed forms of the theorem disagreed by $\sim$1\%. This motivated the convention change above.
\end{verification}

%======================================================================
\section{Claim 2: sharp identified set (paper Theorem 3.3)}

\begin{claim}[Sharp identified set under $\mu$-bounded exposure]\label{cl:sharp}
Suppose $f\le a$ and the hazard model holds for some $\mu\in[\mu_{\mathrm{lo}},\mu_{\mathrm{hi}}]$. Given $(\omega,w,a,f)$,
\[
q(\mu) = 1-\omega\,\frac{w+\mu(a-f)}{w},
\]
is monotone decreasing in $\mu$, and the identified set for $q$ is exactly
$[\,q(\mu_{\mathrm{hi}}),\,q(\mu_{\mathrm{lo}})\,]\cap[0,1]$.
\end{claim}

\begin{proof}
\emph{Step 1 (the curve).} With hazards $(\lambda,\mu\lambda)$, expected civilian deaths are proportional to $\lambda wN + \mu\lambda(a-f)N$, so the civilian share landing in $W$ is $w/(w+\mu(a-f))$ and
\begin{equation}\label{eq:core}
\omega = (1-q)\,\frac{w}{w+\mu(a-f)}.
\end{equation}
Solving \eqref{eq:core} for $q$ gives $q(\mu)$.

\emph{Step 2 (monotonicity).} $\partial q/\partial\mu = -\omega(a-f)/w \le 0$, with equality iff $a=f$ (no civilian adult men) or $\omega=0$.

\emph{Step 3 (validity: every point of the interval with $q(\mu^\ast)\in[0,1]$ is attainable).} Fix any $\mu^\ast\in[\mu_{\mathrm{lo}},\mu_{\mathrm{hi}}]$ such that $q(\mu^\ast)\in[0,1]$ (values outside $[0,1]$ are not shares and are excluded by the final intersection), and any $\lambda>0$ small enough that all implied class death counts are below class populations. The pair $(\lambda,\mu^\ast\lambda)$ is an admissible hazard configuration whose implied $\omega$ satisfies \eqref{eq:core} with $q=q(\mu^\ast)$. Hence $q(\mu^\ast)$ is consistent with the data $(\omega,w,a,f)$: the identified set contains $\{q(\mu):\mu\in[\mu_{\mathrm{lo}},\mu_{\mathrm{hi}}]\}\cap[0,1]$, which by continuity and monotonicity is $[q(\mu_{\mathrm{hi}}),q(\mu_{\mathrm{lo}})]\cap[0,1]$.

\emph{Step 4 (sharpness: nothing outside is attainable).} Conversely, any admissible configuration has some $\mu\in[\mu_{\mathrm{lo}},\mu_{\mathrm{hi}}]$, and \eqref{eq:core} then forces $q=q(\mu)$, which lies in the interval. The intersection with $[0,1]$ imposes the logical range of a share.
\end{proof}

\begin{verification}
\PASS. Sympy confirms the inversion and the derivative. 2{,}000 random parameter draws confirm endpoint ordering. The forward model's three death-shares (combatant, civilian-$\mathrm{AM}$, $W$) sum to 1 symbolically. \emph{Caveat worth knowing when checking manually:} the theorem treats $\omega$ as an exact population quantity; the ``identified set'' is a set for the estimand given error-free inputs. Sampling error enters separately via Claim~\ref{cl:delta} and systematic error via the anchor sensitivity analysis. This is stated in the paper's Discussion (limitation 3).
\end{verification}

%======================================================================
\section{Claim 3: combatants outside AM (paper Remark 3)}

\begin{claim}[$\eta$-slack]\label{cl:remark}
If a fraction $\eta\in[0,\bar\eta]$, $\bar\eta=0.03$, of combatant deaths falls in class $W$, then with $S:=w/(w+\mu(a-f))$ the observable identity becomes $\omega=(1-q)S+\eta q$, so
\[
q_{\mathrm{true}} = \frac{S-\omega}{S-\eta},
\qquad
q_{\mathrm{true}}-q_{\eta=0} \;=\; \frac{\eta(S-\omega)}{S(S-\eta)}
\;=\; \frac{\eta\,q_{\mathrm{true}}}{S}
\;=\; \frac{\eta\,q_{\eta=0}}{S-\eta} \;\ge\; 0 .
\]
Ignoring the slack \emph{understates} $q$, by at most $\bar\eta q_{\eta=0}/(S-\bar\eta)\le 1.4$ percentage points at the Gaza values ($q_{\eta=0}\le 0.26$ and $S\ge 0.597$ for $\mu\le 2$; at $\mu=1$, $S=0.748$ and the bias is $\le 1.1$ points).
\end{claim}

\begin{proof}
Deaths in $W$ now come from two sources: civilian deaths landing in $W$, which is $(1-q)S$ of all deaths as in Claim~\ref{cl:sharp}, plus the fraction $\eta$ of combatant deaths, contributing $\eta q$. Setting $\omega=(1-q)S+\eta q=S-q(S-\eta)$ and solving for $q$ gives $q_{\mathrm{true}}=(S-\omega)/(S-\eta)$. Subtracting $q_{\eta=0}=(S-\omega)/S$,
\[
q_{\mathrm{true}}-q_{\eta=0}
=(S-\omega)\Bigl(\frac{1}{S-\eta}-\frac{1}{S}\Bigr)
=\frac{\eta(S-\omega)}{S(S-\eta)},
\]
which is nonnegative because $\omega\le S$ whenever $q\ge 0$. Substituting $S-\omega=q_{\mathrm{true}}(S-\eta)$ gives the second form $\eta q_{\mathrm{true}}/S$; substituting $S-\omega=q_{\eta=0}S$ gives the third form $\eta q_{\eta=0}/(S-\eta)$. Note the distinction: the bias equals $\eta q/( S-\eta)$ with $q$ the \emph{no-slack} estimate, and $\eta q/S$ with $q$ the \emph{true} share --- the two coincide only to first order in $\eta$.
\end{proof}

\begin{verification}
\FIXED\ (twice). The manuscript originally asserted the shift is ``$-\bar\eta$'' (i.e.\ subtract $\eta$ from $q$). Symbolic verification shows this is wrong in \emph{both} magnitude and sign: the true bias is $+\eta q_{\eta=0}/(S-\eta)$ (understatement, not overstatement, of $q$), and its magnitude is ${\approx}0.011$ at Gaza values ($\mu=1$), not $\eta=0.03$. A second error was caught by external review of a previous version of \emph{this dossier}: the bias chain contained a spurious factor ($\eta q_{\mathrm{true}}/S\cdot S/(S-\eta)$, which double-counts the $(S-\eta)^{-1}$ correction), and the claimed support ``$S\ge 0.70$'' fails for $\mu>1.15$ at Gaza demographics ($S(1.5)=0.664$, $S(2)=0.597$). Both are corrected above; the bias bound is now stated per-$\mu$. Sympy check: all three displayed forms are now symbolically equal to $q_{\mathrm{true}}-q_{\eta=0}$. Direction note unchanged: the slack makes the paper's \emph{upper} bounds conservative in the claim-favourable direction, so no downstream conclusion changes.
\end{verification}

%======================================================================
\section{Claim 4: manpower bound (paper Corollary 3.4)}

\begin{claim}
$q \le M/D$; whenever $\max\{0,q(\mu_{\mathrm{hi}})\}\le\min\{1,\,q(\mu_{\mathrm{lo}}),\,M/D\}$, the joint identified set is $\bigl[\max\{0,q(\mu_{\mathrm{hi}})\},\ \min\{1,\,q(\mu_{\mathrm{lo}}),\,M/D\}\bigr]$; otherwise it is empty and the inputs are jointly infeasible.
\end{claim}

\begin{proof}
Combatants killed cannot exceed combatants existing: $D_M\le M$. Divide by $D>0$. The joint set is the intersection of two intervals, hence the endpoint formulas. (If the intersection is empty, the inputs are jointly infeasible --- this is exactly the event detected by a positive contradiction radius.)
\end{proof}

\begin{verification}
\PASS\ (algebraic identity). One caveat for manual review: $M$ here is the stock at $t_0$; if recruitment during the war is material, $M$ should be the stock-plus-inflow, which \emph{raises} the bound. In the Gaza application the bound is slack by a factor $>2$ ($M/D=0.64$ vs.\ $q\le 0.26$), so recruitment cannot change any conclusion.
\end{verification}

%======================================================================
\section{Claim 5: delta-method standard error (paper \S3)}

\begin{claim}\label{cl:delta}
For $\hat q = 1-\hat\omega(1-\hat f)/\hat w$ with independent inputs,
\[
\widehat{\mathrm{Var}}(\hat q)\approx
\frac{\hat\omega^2}{\hat w^2}\widehat{\mathrm{Var}}(\hat f)
+\frac{(1-\hat f)^2}{\hat w^2}\widehat{\mathrm{Var}}(\hat\omega)
+\frac{\hat\omega^2(1-\hat f)^2}{\hat w^4}\widehat{\mathrm{Var}}(\hat w).
\]
\end{claim}

\begin{proof}
First-order Taylor expansion of $g(\omega,f,w)=1-\omega(1-f)/w$:
\[
\partial_f g = \frac{\omega}{w},\qquad
\partial_\omega g = -\frac{1-f}{w},\qquad
\partial_w g = \frac{\omega(1-f)}{w^2}.
\]
$\mathrm{Var}(\hat q)\approx\sum(\partial g)^2\mathrm{Var}$, using independence (no covariance terms). Squaring the gradients gives the display.
\end{proof}

\begin{verification}
\PASS. All three partial derivatives and the assembled formula verified symbolically. Assumption to flag for manual review: independence of $(\hat\omega,\hat f,\hat w)$ is plausible here because they come from different institutions (deaths records, military-balance estimates, census), but any common-mode error (e.g.\ both $\omega$ and $w$ derived from the same census frame) would add covariance terms. The paper's conservative $\sigma_\omega=0.01$ (5$\times$ the binomial SE) partially absorbs this.
\end{verification}

%======================================================================
\section{Claim 6: precision-weighted aggregation (paper Proposition 2.1)}

\begin{claim}
For independent unbiased sources with variances $\sigma_i^2$ treated as known (or fixed externally; with estimated $\hat\sigma_i^2$ the optimality is conditional on the variance estimates, exact asymptotically), the estimator
\[
\hat\theta_{\mathrm{agg}}=\frac{\sum_i\hat\theta_i/\hat\sigma_i^2}{\sum_i 1/\hat\sigma_i^2}
\]
is minimum-variance among linear unbiased combinations with weights fixed given the variances. With per-source biases $b_i$, the aggregate bias is the precision-weighted mean of the $b_i$. Under $\varepsilon_i$-contamination in which the contaminated report is displaced from the truth by at most $R_i$ (equivalently, the contaminating support is contained in a set of diameter $R_i$ containing the truth), the worst-case aggregate bias is at most $\sum_i \varepsilon_i R_i\,\hat\sigma_{\mathrm{agg}}^2/\hat\sigma_i^2$.
\end{claim}

\begin{proof}
\emph{(i)} Minimise $\sum c_i^2\hat\sigma_i^2$ subject to $\sum c_i=1$. The Lagrangian condition gives $2c_i\hat\sigma_i^2=\kappa$, i.e.\ $c_i\propto 1/\hat\sigma_i^2$; the constraint fixes the normalisation. The objective is strictly convex, so this is the unique minimum (Gauss--Markov in one dimension).

\emph{(ii)} $\E\hat\theta_{\mathrm{agg}}-\theta=\sum_i c_i b_i$ with $c_i=\hat\sigma_{\mathrm{agg}}^2/\hat\sigma_i^2$; direct substitution.

\emph{(iii)} Contamination replaces source $i$'s report with probability $\varepsilon_i$ by a value displaced by at most $R_i$ (diameter of the contaminating support). The expected displacement of source $i$'s contribution is at most $\varepsilon_iR_i$, weighted into the aggregate by $c_i$; summing over $i$ gives the bound.
\end{proof}

\begin{verification}
\PASS. Optimality confirmed on 300 random configurations against random competitor weights; bias formula confirmed symbolically. Nuance for manual review: in (iii) ``diameter of the support'' means the contaminated value cannot leave a set of diameter $R_i$ around truth; if the adversary is unbounded, so is the worst case --- this is why the paper pairs the proposition with the identified-set diameters $R_i$ (removal diameters), which are finite by construction.
\end{verification}

%======================================================================
\section{Claim 7: feasibility and contradiction (paper Theorem 4.3)}

\begin{claim}
Fix $(\hat a,\hat f,\hat D)$ and $[\mu_{\mathrm{lo}},\mu_{\mathrm{hi}}]$. Define $\mathcal F(x)$ as in the paper (eq.~4) and
\[
\rho^{(i)}=\inf_{\omega',w',M'}\max\!\Bigl\{
\tfrac{|\omega'-\hat\omega|}{\sigma_\omega},
\tfrac{|w'-\hat w|}{\sigma_w},
\tfrac{(M'-\hat M)_+}{\sigma_M}\Bigr\}
\ \text{s.t.}\ q^{(i)}\in\mathcal F((\omega',w',\hat a,\hat f,M',\hat D)).
\]
Then (i) $\rho^{(i)}=0$ iff $q^{(i)}\in\mathcal F(\hat x)$, and (ii) if $\rho^{(i)}>0$, every rationalisation $(\omega',w',M')$ has at least one of $|\omega'-\hat\omega|/\sigma_\omega$, $|w'-\hat w|/\sigma_w$, $(M'-\hat M)_+/\sigma_M$ at least $\rho^{(i)}$; that is, at least one input must move by $\rho^{(i)}$ standard errors \emph{in a penalised direction} (downward moves of $M$ are free by construction, since lowering the manpower ceiling can only hurt the claim).
\end{claim}

\begin{proof}
\emph{(i, $\Leftarrow$)} If $q^{(i)}\in\mathcal F(\hat x)$, take $(\omega',w',M')=(\hat\omega,\hat w,\hat M)$: objective $0$.

\emph{(i, $\Rightarrow$)} Suppose $\rho^{(i)}=0$. Take a sequence $(\omega'_n,w'_n,M'_n)$ of feasible points with objective $\to 0$. Because the $M$-penalty is \emph{one-sided}, this does \emph{not} imply $M'_n\to\hat M$; it implies only
\[
\omega'_n\to\hat\omega,\qquad w'_n\to\hat w,\qquad (M'_n-\hat M)_+\to 0,
\]
i.e.\ $\limsup_n M'_n\le\hat M$ (with $M'_n$ possibly far below $\hat M$). We therefore argue via the constraint rather than via convergence of $M'_n$. Each feasible point satisfies (a) $q^{(i)}=1-\omega'_n(w'_n+\mu_n(a-f))/w'_n$ for some $\mu_n\in[\mu_{\mathrm{lo}},\mu_{\mathrm{hi}}]$, and (b) $q^{(i)}D\le M'_n$. For (a): pass to a subsequence with $\mu_n\to\mu^\ast$ in the compact interval; continuity of $(\omega',w',\mu)\mapsto 1-\omega'(w'+\mu(a-f))/w'$ (on $w'>0$) gives $q^{(i)}=1-\hat\omega(\hat w+\mu^\ast(a-f))/\hat w$, so the exposure constraint holds at the observed $(\hat\omega,\hat w)$. For (b): $q^{(i)}D\le M'_n\le \hat M+(M'_n-\hat M)_+\to\hat M$, so $q^{(i)}D\le\hat M$. Hence $q^{(i)}\in\mathcal F(\hat x)$. (Note the repaired step: the naive claim ``objective $\to0$ implies the full vector converges to $(\hat\omega,\hat w,\hat M)$'' is false under the one-sided penalty; only the inequality $q^{(i)}D\le\hat M$ survives the limit, and that is all the theorem needs.)

\emph{(ii)} Let $(\omega',w',M')$ be any point of $C$ (``rationalisation''). By definition of the infimum, $\max\{|\omega'-\hat\omega|/\sigma_\omega,\,|w'-\hat w|/\sigma_w,\,(M'-\hat M)_+/\sigma_M\}\ge\rho^{(i)}$, so at least one of the three \emph{penalised} displacements is $\ge\rho^{(i)}$ SEs. This does not assert ``$M$ moved'' generically: lowering $M'$ costs nothing under the objective (and only shrinks the feasible set), so a positive radius can only ever be certified by an upward move of $M$, or a move of $\omega$ or $w$ in either direction.
\end{proof}

\begin{verification}
\FIXED\ (proof only; statement repaired in wording). External review of a previous version of this dossier correctly observed that the original proof of (i,$\Rightarrow$) claimed $(\omega'_n,w'_n,M'_n)\to(\hat\omega,\hat w,\hat M)$, which is false under the one-sided $M$-penalty (downward moves of $M'$ are free). The repaired proof above passes to the limit through the constraint $q^{(i)}D\le M'_n\le\hat M+o(1)$ instead, and part (ii) now states explicitly which displacements are penalised. The theorem's \emph{conclusion} is unchanged. Numerically, 2{,}000 random infeasible claims all fail to become feasible under $10^{-9}$-perturbations (no boundary pathologies found). The reduction used in the code (radius on the $\omega$-axis $=$ distance from $\hat\omega$ to the interval $[\omega_{\mathrm{needed}}(\mu_{\mathrm{hi}}),\omega_{\mathrm{needed}}(\mu_{\mathrm{lo}})]$) is justified by $\omega_{\mathrm{needed}}$ decreasing in $\mu$ (verified symbolically). \emph{Related error found and fixed elsewhere}: the original \emph{application} of this definition in Section 6 evaluated the radius at the least favourable $\mu$ instead of minimising over the admissible range, overstating $\rho$ ($\gtrsim 20$ claimed; correct values 8--31 depending on anchor and exposure regime, and $\rho=0$ for the 17k endpoint under agnostic exposure on the MoH anchor). The definition and theorem themselves were always correct.
\end{verification}

%======================================================================
\section{Claim 8: bias-robust sensitivity band (paper Proposition 5.1, restated)}

\begin{claim}
Let $I_0=[L_0,U_0]$ be the $1-\alpha$ credible interval under face-value sources, posterior supported on the identified set implied by the sources used. Each source $i$ is independently contaminated with probability $\varepsilon_i$; removing the sources in $T$ leaves an identified set of diameter $R_T$ ($R_i:=R_{\{i\}}$). Then: (i) conditional on realised pattern $T$, every attainable $1-\alpha$ credible interval lies within $[L_0-R_T,\,U_0+R_T]$; (ii) the \emph{expected} endpoint displacement is at most $\sum_i\varepsilon_iR_i$ plus $O(\varepsilon^2)$ terms involving multi-removal diameters, so $I_\varepsilon=[L_0-\sum_i\varepsilon_iR_i,\ U_0+\sum_i\varepsilon_iR_i]$ is a first-order sensitivity band, \emph{not} a uniform envelope over realised intervals; (iii) if the posterior mixture weight on contaminated components is at most $\tau:=1-\prod_i(1-\varepsilon_i)$ and the face-value posterior density exceeds $m>0$ on the segment between the clean and contaminated quantiles,
\[
|\hat Q_\varepsilon(u)-\hat Q_0(u)|\;\le\; m^{-1}\tau,
\qquad u\in\{\alpha/2,\,1-\alpha/2\}.
\]
\end{claim}

\begin{proof}
\emph{(i)} Condition on $T$: honest sources outside $T$ confine $\theta$ to the removal identified set of diameter $R_T$, which contains the face-value value; any posterior formed under adversarial reports from $T$ is supported there, so each credible-interval endpoint is within $R_T$ of the corresponding face-value endpoint.

\emph{(ii)} The pattern probability is $P(T)=\prod_{i\in T}\varepsilon_i\prod_{i\notin T}(1-\varepsilon_i)$, so the expected displacement is at most $\sum_T P(T)R_T = \sum_i\varepsilon_iR_i\prod_{j\ne i}(1-\varepsilon_j) + \sum_{|T|\ge2}P(T)R_T \le \sum_i\varepsilon_iR_i + O(\varepsilon^2)$. This is a statement about the \emph{mean} over the contamination pattern. A quantile (interval endpoint) of the realised posterior can move by up to $R_T$ on a contamination event of probability $\varepsilon_i$: with one source, $\varepsilon=0.04$, $R=100$, a contaminated component at distance $\approx R$ carrying $4\%$ posterior mass moves the $97.5\%$ quantile by $\approx R$, far beyond $\varepsilon R=4$. Hence the band is first-order in $\varepsilon$, exact when at most one source can be contaminated \emph{and} displacement is measured in expectation.

\emph{(iii)} If the posterior mixture weight on contaminated components is $\le\tau$, the contaminated and face-value posterior CDFs differ by at most $\tau$ uniformly; the mean value theorem under the density floor $m$ (on the whole segment between the two quantiles) gives the display. The weight condition does \emph{not} follow from the prior probabilities alone: posterior component weights are $\propto P(T)\,Z_T$ where $Z_T$ is the component marginal likelihood, so an adversarial component that fits the data better than the clean model carries posterior weight exceeding $\tau$.
\end{proof}

\begin{verification}
\FIXED\ (statement was materially over-claimed). External review of a previous version of this dossier exhibited the counterexample reproduced in (ii): the original statement asserted that $I_\varepsilon$ \emph{contains every} attainable credible interval, which is false for tail quantiles --- the mixture argument controls an expectation, not an envelope. Numerically confirmed: with $\varepsilon=0.04$, $R=100$, clean posterior $\mathcal N(0,0.1^2)$ and contaminated component $\mathcal N(100,0.1^2)$, the pooled $97.5\%$ quantile is $99.97$ while $\varepsilon R=4$. The same review identified the posterior-weight gap in the old part (ii): mixture weights after conditioning are $P(T)Z_T$-proportional, not $P(T)$. The proposition has been restated in the manuscript as a \emph{patternwise envelope} $+$ \emph{first-order sensitivity band} $+$ \emph{quantile bound under an explicit posterior-weight condition}, and all prose referring to it as a ``credible interval that contains every adversarial interval'' has been reworded. Downstream use in Section 6 is unaffected in substance: the quoted $25.1\%\to 28.7\%$ inflation is the first-order band $I_\varepsilon$, now labelled as such; the paper's headline rejections rest on the bounds and the contradiction radius, not on this proposition.
\end{verification}

%======================================================================
\section{Section 6 arithmetic (spot summary)}

All application numbers are generated by \texttt{analysis/validate\_bounds.py}. Key identities re-verified here at $(w,a,f,\omega)=(0.733,0.267,0.020,0.560)$:
\begin{center}
\begin{tabular}{@{}lll@{}}
\toprule
Quantity & Value & Check \\
\midrule
$q(1)$ & $25.13\%$ & symbolic $+$ numeric \\
$q(1.5)$ & $15.69\%$ & numeric (prior table misprinted $15.72\%$) \\
$q(2)$ & $6.26\%$ & numeric \\
root of $q(\mu)=0$ & $\mu^\ast=2.33$ & numeric; matches the $q{=}0$ symmetry discussion \\
$\mu$ needed for $q=17/70$ & $1.04$ & inversion; $\ge 1$ so arithmetically feasible \\
$\mu$ needed for $q=25/70$ & $0.44$ & inversion; $<1$ so infeasible for any admissible $\mu$ \\
$M/D$ & $0.643$ & non-binding vs.\ all bounds above \\
\bottomrule
\end{tabular}
\end{center}

%======================================================================
\section{Errors found by this verification (all fixed in the manuscript)}

\begin{enumerate}[leftmargin=*]
\item \textbf{Remark 3 ($\eta$-slack) was wrong in sign and magnitude.} Claimed shift ``$-\bar\eta$''; correct bias is $+\eta q/(S-\eta)\approx 1$ percentage point at Gaza values (ignoring the slack \emph{understates} $q$). Rewritten with the exact expression. Downstream conclusions unaffected (direction is claim-favourable).
\item \textbf{The contradiction radius was originally computed at the wrong $\mu$.} Definition 4.2 grants the claim the most favourable $\mu$ in the admissible range; the original Section 6 evaluated at the least favourable, producing ``$\rho\gtrsim 20$''. Corrected values: $\rho_\omega\approx 7.9$--$31$ depending on anchor and exposure regime, and $\rho=0$ for the 17k endpoint under agnostic exposure on the MoH anchor (it is rejected by exposure calibration, not arithmetic). Disclosed in the response letter.
\item \textbf{Demographic classes did not partition the population.} $a=0.255$ (males 18--60) left men over 60 unassigned while proofs used $a=1-w$. Convention changed to $\mathrm{AM}=$ males $18+$, $a=0.267$; all numbers regenerated.
\end{enumerate}

\section{Errors found by external review of this dossier (second round, all fixed)}

A previous version of this dossier was independently reviewed; the following defects in the \emph{dossier and manuscript statements} (as distinct from the application numbers) were confirmed and corrected:

\begin{enumerate}[leftmargin=*]
\item \textbf{Claim 3's displayed bias chain contained a spurious factor.} The chain $\eta q_{\mathrm{true}}/S\cdot S/(S-\eta)$ double-counts the $(S-\eta)^{-1}$ correction; the correct equalities are $\eta(S-\omega)/(S(S-\eta))=\eta q_{\mathrm{true}}/S=\eta q_{\eta=0}/(S-\eta)$. Verified symbolically; fixed here and in manuscript Remark~3, which now names which $q$ appears in which form.
\item \textbf{Claim 3's numerical support ``$S\ge 0.70$'' was false for $\mu>1.15$.} At Gaza demographics $S(1.5)=0.664$, $S(2)=0.597$. The bias bound is now stated per-$\mu$ (max $1.4$ points at $\mu=2$).
\item \textbf{Claim 7's proof mishandled the one-sided $M$-penalty.} ``Objective $\to0$ implies the full vector converges'' is false when downward moves of $M$ are free. Repaired by passing to the limit through the constraint $q^{(i)}D\le M'_n\le\hat M+o(1)$; theorem conclusion unchanged. Part (ii) now names the penalised displacements explicitly. Manuscript Theorem~4.3 statement and appendix proof updated.
\item \textbf{Claim 8 (Proposition 5.1) was materially over-claimed as a credible-interval envelope.} A one-source counterexample ($\varepsilon=0.04$, $R=100$) moves the $97.5\%$ quantile by ${\approx}100\gg\varepsilon R=4$. The mixture argument controls an \emph{expectation}; additionally, posterior mixture weights are $P(T)Z_T$-proportional, not $P(T)$. Restated in the manuscript as patternwise envelope $+$ first-order sensitivity band $+$ quantile bound under an explicit posterior-weight condition; all ``contains every adversarial interval'' prose removed. The quoted Gaza inflation $25.1\%\to28.7\%$ is unchanged as a first-order band; no headline conclusion rests on this proposition.
\item \textbf{Smaller repairs.} Claim 2 Step 3 now restricts attainability to $q(\mu)\in[0,1]$; Claim 4 states the nonempty-intersection condition explicitly; Claim 6 now says variances are treated as known and that $R_i$ bounds displacement \emph{from the truth}; the arithmetic table's $q(1.5)$ was misprinted as $15.72\%$ (correct: $15.69\%$; the manuscript's rounded $15.7\%$ was already correct).
\end{enumerate}

\end{document}
